\section{Schemes}
\label{section-schemes}

\begin{definition}
\label{definition-scheme}
\begin{history}
In \cite{EGA1} what we call a scheme was called a ``pre-sch\'ema'' and the
name ``sch\'ema'' was reserved for what is a separated scheme in the
Stacks project. In the second edition \cite{EGA1-second} the terminology
was changed to the terminology that is now standard. However, one may
occasionally encounter the terminology ``prescheme'', for example in
\cite{Murre-lectures}.
\end{history}
A {\it scheme} is a locally ringed space with the property that
every point has an open neighbourhood which is an affine scheme.
A {\it morphism of schemes} is a morphism of locally
ringed spaces. The category of schemes will be denoted
$\Sch$.
\end{definition}

\noindent
Let $X$ be a scheme.
We will use the following (very slight) abuse of language.
We will say $U \subset X$ is an {\it affine open}, or an {\it open affine}
if the open subspace $U$ is an affine scheme. We will often
write $U = \Spec(R)$ to indicate that $U$ is isomorphic
to $\Spec(R)$ and moreover that we will identify (temporarily)
$U$ and $\Spec(R)$.

\begin{lemma}
\label{lemma-open-subspace-scheme}
Let $X$ be a scheme. Let $j : U \to X$ be an open immersion
of locally ringed spaces. Then $U$ is a scheme. In particular,
any open subspace of $X$ is a scheme.
\end{lemma}

\begin{proof}
Let $U \subset X$. Let $u \in U$.
Pick an affine open neighbourhood $u \in V \subset X$.
Because standard opens of $V$ form a basis of the topology
on $V$ we see that there exists a $f\in \mathcal{O}_V(V)$
such that $u \in D(f) \subset U$. And $D(f)$ is an affine scheme
by Lemma \ref{lemma-standard-open-affine}. This proves that every point
of $U$ has an open neighbourhood which is affine.
\end{proof}

\noindent
Clearly the lemma (or its proof) shows that any scheme
$X$ has a basis (see Topology, Section \ref{topology-section-bases})
for the topology consisting of affine opens.

\begin{example}
\label{example-not-affine}
Let $k$ be a field.
An example of a scheme which is not affine is
given by the open subspace
$U = \Spec(k[x, y]) \setminus \{ (x, y)\}$
of the affine scheme $X =\Spec(k[x, y])$.
It is covered by two affines, namely $D(x) = \Spec(k[x, y, 1/x])$
and $D(y) = \Spec(k[x, y, 1/y])$ whose intersection is
$D(xy) = \Spec(k[x, y, 1/xy])$. By the sheaf property
for $\mathcal{O}_U$ there is an exact sequence
$$
0 \to
\Gamma(U, \mathcal{O}_U) \to
k[x, y, 1/x] \times k[x, y, 1/y] \to
k[x, y, 1/xy]
$$
We conclude that the map $k[x, y] \to \Gamma(U, \mathcal{O}_U)$
(coming from the morphism $U \to X$) is an isomorphism.
Therefore $U$ cannot be affine since if it was then by
Lemma \ref{lemma-category-affine-schemes} we would have $U \cong X$.
\end{example}
